\section{Conclusion}
We gave an overview of prompt injection attack techniques on LLMs, beginning with the formal definition of prompt injection and the basic threat model shared amongst attack approaches.
Besides the temporal development process of prompt injection approaches, we structured the techniques into two distinct classes.
One consisting of word-level attack prompt construction and the other containing token-level gradient-based attack prompt optimization.\\

Following the development of prompt injection attacks, we presented a manual approach to creating attack prompts based on general prompting guidelines.
We then examined the next development step by discussing semi-automated approaches.
By comparing the data of several case studies, we concluded that better instruction following, enabled by larger models and better fine-tuning, makes LLMs more susceptible to prompt injection attacks.
At the same time, the case studies indicated that RLHF is empirically more robust against prompt injection attacks than supervised fine-tuning.\\

We proceeded to present the second class of gradient-based approaches, due to the observation that older attack techniques become less effective against newer LLMs.
Comparing gradient-based approaches to LLM-based attack prompt rephrasing, we highlighted that gradient optimization has the advantages of high attack effectiveness and easy specialization via multiple loss terms.
We also pointed out the higher complexity of gradient-based approaches, which is due to their need for a shadow model and higher computational effort.
In reference to the development progress of prompt injection attacks, we noted the more realistic attack scenarios discussed in newer studies.\\

Lastly, we discussed the fundamental problem of instruction-data integration that enables the ongoing high success rates of prompt injection attacks on state-of-the-art LLMs.
We presented three types of defense strategies.
The band-aid solutions that are currently widely adopted, deeper architectural adjustments with potentially better generalizability, and interpretability-based defenses that could offer a conceptual solution to prompt injection attacks if they become applicable at scale.\\

Overall, prompt injection attacks against LLMs remain one of the most dominant threats to LLM-based applications.
Considering the increasing adoption and impact of such applications, novel defense mechanisms like interpretability-based defenses are needed to underlay further efforts of LLM-based automation.